\documentclass{amsart}
% For dealing with references we use the comment environment
\usepackage{verbatim}
\newenvironment{reference}{\comment}{\endcomment}
%\newenvironment{reference}{}{}
\newenvironment{slogan}{\comment}{\endcomment}
\newenvironment{history}{\comment}{\endcomment}

% For commutative diagrams we use Xy-pic
\usepackage[all]{xy}

% We use 2cell for 2-commutative diagrams.
\xyoption{2cell}
\UseAllTwocells

% We use multicol for the list of chapters between chapters
\usepackage{multicol}

% This is generally recommended for better output
\usepackage{lmodern}
\usepackage[T1]{fontenc}

% For cross-file-references

% Package for hypertext links:
\usepackage{hyperref}

% For any local file, say "hello.tex" you want to link to please

% Theorem environments.
%
\theoremstyle{plain}
\newtheorem{theorem}[subsection]{Theorem}
\newtheorem{proposition}[subsection]{Proposition}
\newtheorem{lemma}[subsection]{Lemma}

\theoremstyle{definition}
\newtheorem{definition}[subsection]{Definition}
\newtheorem{example}[subsection]{Example}
\newtheorem{exercise}[subsection]{Exercise}
\newtheorem{situation}[subsection]{Situation}

\theoremstyle{remark}
\newtheorem{remark}[subsection]{Remark}
\newtheorem{remarks}[subsection]{Remarks}

\numberwithin{equation}{subsection}

% Macros
%
\def\lim{\mathop{\mathrm{lim}}\nolimits}
\def\colim{\mathop{\mathrm{colim}}\nolimits}
\def\Spec{\mathop{\mathrm{Spec}}}
\def\Hom{\mathop{\mathrm{Hom}}\nolimits}
\def\Ext{\mathop{\mathrm{Ext}}\nolimits}
\def\SheafHom{\mathop{\mathcal{H}\!\mathit{om}}\nolimits}
\def\SheafExt{\mathop{\mathcal{E}\!\mathit{xt}}\nolimits}
\def\Sch{\mathit{Sch}}
\def\Mor{\mathop{\mathrm{Mor}}\nolimits}
\def\Ob{\mathop{\mathrm{Ob}}\nolimits}
\def\Sh{\mathop{\mathit{Sh}}\nolimits}
\def\NL{\mathop{N\!L}\nolimits}
\def\CH{\mathop{\mathrm{CH}}\nolimits}
\def\proetale{{pro\text{-}\acute{e}tale}}
\def\etale{{\acute{e}tale}}
\def\QCoh{\mathit{QCoh}}
\def\Ker{\mathop{\mathrm{Ker}}}
\def\Im{\mathop{\mathrm{Im}}}
\def\Coker{\mathop{\mathrm{Coker}}}
\def\Coim{\mathop{\mathrm{Coim}}}

% Boxtimes
%
\DeclareMathSymbol{\boxtimes}{\mathbin}{AMSa}{"02}

%
% Macros for moduli stacks/spaces
%
\def\QCohstack{\mathcal{QC}\!\mathit{oh}}
\def\Cohstack{\mathcal{C}\!\mathit{oh}}
\def\Spacesstack{\mathcal{S}\!\mathit{paces}}
\def\Quotfunctor{\mathrm{Quot}}
\def\Hilbfunctor{\mathrm{Hilb}}
\def\Curvesstack{\mathcal{C}\!\mathit{urves}}
\def\Polarizedstack{\mathcal{P}\!\mathit{olarized}}
\def\Complexesstack{\mathcal{C}\!\mathit{omplexes}}
% \Pic is the operator that assigns to X its picard group, usage \Pic(X)
% \Picardstack_{X/B} denotes the Picard stack of X over B
% \Picardfunctor_{X/B} denotes the Picard functor of X over B
\def\Pic{\mathop{\mathrm{Pic}}\nolimits}
\def\Picardstack{\mathcal{P}\!\mathit{ic}}
\def\Picardfunctor{\mathrm{Pic}}
\def\Deformationcategory{\mathcal{D}\!\mathit{ef}}

\setlength{\emergencystretch}{3em}
\begin{document}
\title[Stacks Project, Tag 0EA2]{325 --- Effective exceptional canonical divisors characterize rational surface singularities}
\maketitle
\noindent Copyright (C) 2005--2025 Johan de Jong. Stacks Project authors. Source: \href{https://stacks.math.columbia.edu/tag/0EA2}{Tag 0EA2}.

\noindent Editorial difficulty note (not author text). Difficulty H2: The proof must convert an effective exceptional canonical divisor into equality of derived pushforwards. Trace and Grauert-Riemenschneider vanishing are combined with a depth-two isomorphism argument, and duality then gives rationality; minimality additionally forces the exceptional divisor to vanish..

\noindent Editorial provenance note (not author text). History: excerpt from revision \texttt{a0444\allowbreak 6e57e\allowbreak c1fbc\allowbreak 252a8\allowbreak 71afc\allowbreak ec775\allowbreak 2fb28\allowbreak 07b14}, pione.tex, lines 6167--6217. Reference presentation was adapted; original-author context and core support blocks were retained or added. The mathematical wording is unchanged. Licensed under GNU FDL 1.2 or later; no invariant sections or cover texts. See ../licenses/COPYING. Context blocks reproduce author definitions and supporting results; they are not additional counted problems.

\begin{lemma}
\label{lemma-characterize-rational-singularity}
Let $A$ be a Noetherian normal local domain of dimension $2$.
Assume $A$ is Nagata, has a dualizing module $\omega_A$, and has a
resolution of singularities $f : X \to \Spec(A)$.
Let $\omega_X$ be as in Resolution of Surfaces,
Remark \href{https://stacks.math.columbia.edu/tag/0B4R}{remark dualizing setup (Tag 0B4R)}.
If $\omega_X \cong \mathcal{O}_X(E)$ for some effective
Cartier divisor $E \subset X$ supported on the exceptional
fibre, then $A$ defines a rational singularity.
If $f$ is a minimal resolution, then $E = 0$.
\end{lemma}

\begin{proof}
There is a trace map $Rf_*\omega_X \to \omega_A$, see
Duality for Schemes, Section \href{https://stacks.math.columbia.edu/tag/0AWG}{section trace (Tag 0AWG)}.
By Grauert-Riemenschneider
(Resolution of Surfaces,
Proposition \href{https://stacks.math.columbia.edu/tag/0AXD}{proposition Grauert Riemenschneider (Tag 0AXD)})
we have $R^1f_*\omega_X = 0$.
Thus the trace map is a map $f_*\omega_X \to \omega_A$.
Then we can consider
$$
\mathcal{O}_{\Spec(A)} = f_*\mathcal{O}_X \to f_*\omega_X \to \omega_A
$$
where the first map comes from the map
$\mathcal{O}_X \to \mathcal{O}_X(E) = \omega_X$ which is
assumed to exist in the statement of the lemma.
The composition is an isomorphism by Divisors, Lemma
\href{https://stacks.math.columbia.edu/tag/0AVM}{lemma check isomorphism via depth and ass (Tag 0AVM)}
as it is an isomorphism over the punctured spectrum of $A$
(by the assumption in the lemma and the fact that $f$ is an isomorphism
over the punctured spectrum) and $A$ and $\omega_A$
are $A$-modules of depth $2$ (by
Algebra, Lemma \href{https://stacks.math.columbia.edu/tag/031S}{lemma criterion normal (Tag 031S)} and
Dualizing Complexes, Lemma \href{https://stacks.math.columbia.edu/tag/0AWE}{lemma depth dualizing module (Tag 0AWE)}).
Hence $f_*\omega_X \to \omega_A$ is surjective whence an isomorphism.
Thus $Rf_*\omega_X = \omega_A$ which by duality implies
$Rf_*\mathcal{O}_X = \mathcal{O}_{\Spec(A)}$.
Whence $H^1(X, \mathcal{O}_X) = 0$ which implies that $A$
defines a rational singularity (see discussion in
Resolution of Surfaces, Section
\href{https://stacks.math.columbia.edu/tag/0AXE}{section bounded (Tag 0AXE)} in particular
Lemmas \href{https://stacks.math.columbia.edu/tag/0B4Q}{lemma regular rational (Tag 0B4Q)} and
\href{https://stacks.math.columbia.edu/tag/0AXF}{lemma exact sequence (Tag 0AXF)}).
If $f$ is minimal, then $E = 0$ because the map
$f^*\omega_A \to \omega_X$ is surjective by
a repeated application of Resolution of Surfaces, Lemma
\href{https://stacks.math.columbia.edu/tag/0B64}{lemma dualizing blow up rational (Tag 0B64)}
and $\omega_A \cong A$ as we've seen above.
\end{proof}
\end{document}
